\chapter{Rice分布}
    \section{定義}
        \begin{shadebox}
            $a, \sigma>0$とする。
            Rice分布の連続確率密度関数を次式で定義する。
            \[ f(x;a,\sigma) := \frac{x^2}{\sigma^2}\exp\left(-\frac{x^2+a^2}{2\sigma^2}\right)I_0\left(\frac{ax}{\sigma^2}\right) \]
            ここに$I_0$は0次の第一種変形Bessel関数である。
        \end{shadebox}
    \section{正規分布する成分を持つ2次元ベクトルと定数ベクトルとの関係}
        \subsection{記号の準備}
            以下で用いる記号ここでを定義しておく。
            $\bm{a} := [a_1,a_2]^\top\in\realNumbers^2,\;a := \norm{\bm{a}}_2$とする。
            $\tilde{\bm{X}}\in\realNumbers^2$の各成分が正規分布$N(0,\sigma^2)$に従う独立な確率変数であるとする。
            $\bm{X} := \bm{a} + \tilde{\bm{X}}$とし、$\bm{X}$の第$i$成分を$X_i$と表す。
        \subsection{ノルムとの関係}
            $\norm{\bm{X}}_2$は母数$\sigma$のRayleigh分布に従う。
            \begin{proof}
                \quad\par
                座標系を適当に回転させて$\bm{a} = [a,0]^\top$とでき、この操作は$\norm{\bm{X}}_2$を変化させない。
                よって一般性を失わずに$\bm{a} = [a,0]^\top$であるとする。
                \par
                $X_i$の確率密度分布を$f_i$とする。
                $f_1$は$N(0,\sigma^2)$の、$f_2$は$N(a,\sigma^2)$の確率密度関数である。
                $r \geq 0$とし、原点を中心とする半径$r$の円盤領域を$\Omega(r)$とする。
                \begin{align*}
                    \Pr{\norm{\bm{X}}_2 \leq r} &= \int_{\Omega(r)} f_1(x_1)f_2(x_2) \mathrm{d}x_1\mathrm{d}x_2 = \integrate{0}{r}{\integrate{0}{2\pi}{f_1(r'\cos\theta)f_2(r'\sin\theta)r'}{}{\theta}}{}{r'} \\
                    &= \integrate{0}{r}{\integrate{0}{2\pi}{\frac{1}{2\pi\sigma^2}\exp\left(-\frac{1}{2\sigma^2}\left((r'\cos\theta-a)^2+r'^2\sin\theta\right)\right)r'}{}{\theta}}{}{r'} \\
                    &= \frac{1}{\sigma^2}\integrate{0}{r}{r'\exp\left(-\frac{1}{2\sigma^2}\left(r'^2+a^2\right)^2\right)\underbrace{\frac{1}{2\pi}\integrate{0}{2\pi}{\exp\left(\frac{ar'\cos\theta}{\sigma^2}\right)}{}{\theta}}_{(A)}}{}{r'}
                \end{align*}
                (A)を評価する。
                \begin{align*}
                    (A) &= \frac{1}{2\pi}\integrate{0}{2\pi}{\exp\left(\frac{ar'\sin\theta}{\sigma^2}\right)}{}{\theta} = \frac{1}{2\pi}\integrate{0}{2\pi}{\exp\left(\frac{-ar'\sin\theta}{\sigma^2}\right)}{}{\theta} \\
                    &\phantom{=} (\text{被積分関数の$2\pi$周期性を用いた}) \\
                    &= J_0\left(i\frac{ar'}{\sigma^2}\right) = I_0\left(\frac{ar'}{\sigma^2}\right)
                \end{align*}
                よって次式が成り立つ。
                \[ \Pr{\norm{\bm{X}}_2 \leq r} = \frac{1}{\sigma^2}\integrate{0}{r}{r'\exp\left(-\frac{1}{2\sigma^2}\left(r'^2+a^2\right)^2\right)I_0\left(\frac{ar'}{\sigma^2}\right)}{}{r'} \]
                $\norm{\bm{X}}_2$の確率密度関数は次式である。
                \[ \partDerivLong{\Pr{\norm{\bm{X}}_2 \leq r}}{r}{} = \frac{1}{\sigma^2}r\exp\left(-\frac{1}{2\sigma^2}\left(r^2+a^2\right)^2\right)I_0\left(\frac{ar}{\sigma^2}\right) \]
            \end{proof}
        \subsection{偏角との関係}
            \begin{shadebox}
                $\theta_0 := \Arg{a_1 + ia_2},\;\Theta := \Arg{X_1 + iX_2}$とする。
                $\Theta$の確率密度関数は次式である。
                {\small \[ \rho(\theta;\bm{a},\sigma) = \frac{1}{2\pi}\exp\left(-\frac{a^2}{2\sigma^2}\right)\left[1+\frac{a\cos(\theta-\theta_0)}{\sigma}\sqrt{\frac{\pi}{2}}\exp\left(-\frac{a^2\cos^2(\theta-\theta_0)}{2\sigma^2}\right)\left(1+\erf\left(\frac{a\cos(\theta-\theta_0)}{\sqrt{2}\sigma}\right)\right)\right] \]}
            \end{shadebox}
            \begin{proof}
                \quad\par
                まず$\bm{a} = [a,0]^\top$の場合を考え、その後で一般化する。
                $X_i$の確率密度分布を$f_i$とする。
                $f_1$は$N(0,\sigma^2)$の、$f_2$は$N(a,\sigma^2)$の確率密度関数である。
                系は$\Theta$に関して偶対称であるから、$\Theta\in[0,\pi]$の範囲で考えた結果を$\Theta<0$の場合に偶対称に拡張できる。
                \[ \Pr{0\leq\Theta\leq\theta} = \integrate{\theta'=0}{\theta}{\integrate{r'=0}{\infty}{\frac{r'}{2\pi\sigma^2}\exp\left(-\frac{1}{2\sigma^2}\left((r'\cos\theta'-a)^2+r'^2\sin^2\theta'\right)\right)}{}{r'}}{}{\theta'} \]
                よって$\Theta$の確率密度関数を$\rho(\Theta;a,\sigma)$とすると
                \begin{align*}
                    \rho(\Theta;a,\sigma) &= \partDerivLong{\Pr{0\leq\Theta\leq\theta}}{\theta}{} = \integrate{r'=0}{\infty}{\frac{r'}{2\pi\sigma^2}\exp\left(-\frac{1}{2\sigma^2}\left((r'\cos\theta'-a)^2+r'^2\sin^2\theta'\right)\right)}{}{r'} \\\
                    &= \frac{1}{2\pi}\exp\left(-\frac{a^2\sin^2\theta}{2\sigma^2}\right)\underbrace{\integrate{r'=0}{\infty}{\frac{r'}{\sigma^2}\exp\left(-\frac{(r'-a\cos\theta)^2}{2\sigma^2}\right)}{}{r'}}_{(A)}
                \end{align*}
                ここで
                \[ -\partDerivLong{\exp\left(-\frac{(r'-a\cos\theta)^2}{2\sigma^2}\right)}{r'}{} = \frac{r'-a\cos\theta}{\sigma^2}\exp\left(-\frac{(r'-a\cos\theta)^2}{2\sigma^2}\right) \]
                を利用し、
                \begin{align*}
                    (A) &= \left[-\exp\left(-\frac{(r'-a\cos\theta)^2}{2\sigma^2}\right)\right]_{r'=0}^\infty + \integrate{r'=0}{\infty}{}{}{}
                \end{align*}
            \end{proof}
